%===============================================================================
% DOCUMENTO LATEX COMPLETO Y AUTÓNOMO
% Apuntes de Derecho Procesal Civil y Comercial
% Unidades 22 y 23: El Juicio Ejecutivo y Cumplimiento de la Sentencia de Remate
%===============================================================================

\documentclass[12pt,oneside]{book}

% CODIFICACIÓN
\usepackage[utf8]{inputenc}
\usepackage[T1]{fontenc}
\usepackage[spanish,es-nodecimaldot]{babel}

% MÁRGENES
\usepackage[
    top=2.5cm,
    bottom=2.5cm,
    left=3cm,
    right=2.5cm,
    headheight=15pt
]{geometry}

% MATEMÁTICA
\usepackage{amsmath}
\usepackage{amssymb}
\usepackage{mathtools}

% IMÁGENES
\usepackage{graphicx}
\usepackage{float}
\usepackage{caption}
\usepackage{subcaption}

% TABLAS
\usepackage{booktabs}
\usepackage{array}
\usepackage{longtable}
\usepackage{tabularx}

% LISTAS
\usepackage{enumitem}

% COLORES Y CAJAS
\usepackage{xcolor}
\usepackage[most]{tcolorbox}

% ENCABEZADOS
\usepackage{fancyhdr}

% HIPERVÍNCULOS Y METADATOS
\usepackage[
    colorlinks=true,
    linkcolor=blue!50!black,
    citecolor=green!40!black,
    urlcolor=blue!60!black,
    pdftitle={Apuntes de Derecho Procesal Civil y Comercial},
    pdfauthor={Isaac Edgar Camacho Ocampo},
    pdfsubject={Apuntes universitarios},
    bookmarks=true
]{hyperref}

%===============================================================================
% CONFIGURACIÓN GENERAL
%===============================================================================

\setlength{\parindent}{0pt}
\setlength{\parskip}{0.5em}

\setlist[itemize]{
    topsep=0.3em,
    itemsep=0.2em
}

\setlist[enumerate]{
    topsep=0.3em,
    itemsep=0.2em
}

\clubpenalty=10000
\widowpenalty=10000

\setcounter{tocdepth}{2}

%===============================================================================
% DEFINICIÓN DE COMANDO PARA DATOS
%===============================================================================

\newcommand{\universidad}{Universidad de Buenos Aires}
\newcommand{\facultad}{Facultad de Derecho}
\newcommand{\materia}{Derecho Procesal Civil y Comercial}
\newcommand{\titulo}{Unidades 22 y 23: El Juicio Ejecutivo y el Cumplimiento de la Sentencia de Remate}
\newcommand{\autor}{Isaac Edgar Camacho Ocampo}
\newcommand{\anio}{2025}

%===============================================================================
% DEFINICIÓN DE CAJAS ACADÉMICAS
%===============================================================================

\newtcolorbox{definition}[1]{
    enhanced,
    breakable,
    title={Definición: #1},
    fonttitle=\bfseries,
    colback=blue!3,
    colframe=blue!50!black,
    boxrule=0.8pt,
    arc=2mm
}

\newtcolorbox{important}{
    enhanced,
    breakable,
    title={Importante},
    fonttitle=\bfseries,
    colback=yellow!8,
    colframe=orange!70!black,
    boxrule=0.8pt,
    arc=2mm
}

\newtcolorbox{example}[1]{
    enhanced,
    breakable,
    title={Ejemplo: #1},
    fonttitle=\bfseries,
    colback=green!4,
    colframe=green!40!black,
    boxrule=0.8pt,
    arc=2mm
}

\newtcolorbox{formula}[1]{
    enhanced,
    breakable,
    title={Fórmula / Regla: #1},
    fonttitle=\bfseries,
    colback=purple!4,
    colframe=purple!50!black,
    boxrule=0.8pt,
    arc=2mm
}

\newtcolorbox{note}{
    enhanced,
    breakable,
    title={Nota},
    fonttitle=\bfseries,
    colback=white,
    colframe=gray!50,
    boxrule=0.6pt,
    arc=2mm
}

%===============================================================================
% ENCABEZADOS Y PIES
%===============================================================================

\pagestyle{fancy}
\fancyhf{}

\fancyhead[L]{\nouppercase{\leftmark}}
\fancyhead[R]{\materia}

\fancyfoot[C]{\thepage}

\renewcommand{\headrulewidth}{0.4pt}
\renewcommand{\footrulewidth}{0pt}

\fancypagestyle{plain}{
    \fancyhf{}
    \fancyfoot[C]{\thepage}
    \renewcommand{\headrulewidth}{0pt}
}

%===============================================================================
% INICIO DEL DOCUMENTO
%===============================================================================

\begin{document}

%===============================================================================
% PORTADA
%===============================================================================

\begin{titlepage}

\thispagestyle{empty}

\begin{center}

\vspace*{1cm}

{\large \textsc{\universidad}}

\vspace{0.2cm}

{\large \textsc{\facultad}}

\vspace{3cm}

{\Huge \textbf{\materia}}

\vspace{1cm}

{\LARGE \titulo}

\vspace{0.8cm}

\textit{El conocimiento de los procedimientos ejecutivos es la llave para entender cómo se materializa la tutela jurisdiccional.}

\vspace{1.2cm}

\rule{0.70\textwidth}{0.5pt}

\vspace{0.5cm}

{\large Apuntes de estudio}

\vspace{0.5cm}

\rule{0.70\textwidth}{0.5pt}

\vfill

{\large \autor}

\vspace{0.3cm}

{\large \anio}

\end{center}

\vfill

\begin{flushleft}
\textit{Este es un modesto aporte a los estudiantes de ``Derecho Procesal Civil y Comercial'' no pretende sustituir el material oficial de la materia.}
\end{flushleft}

\end{titlepage}

%===============================================================================
% ÍNDICE
%===============================================================================

\tableofcontents

%===============================================================================
% INTRODUCCIÓN
%===============================================================================

\chapter*{Introducción}

Estos apuntes corresponden a las Unidades 22 y 23 del programa de Derecho Procesal Civil y Comercial, que abordan dos temáticas fundamentales dentro del proceso de ejecución: el \textbf{juicio ejecutivo} y el \textbf{cumplimiento de la sentencia de remate}.

El juicio ejecutivo constituye un mecanismo procesal diseñado para permitir que el acreedor pueda hacer efectiva su obligación de manera más expeditiva que a través de un proceso de conocimiento ordinario. Requiere un título ejecutivo pero presupone una lógica completamente distinta a la de la ejecución de sentencias.

El cumplimiento de la sentencia de remate representa la etapa coercitiva del proceso: es decir, el momento en que el Estado, a través del poder jurisdiccional, impone coercitivamente su decisión cuando las personas no han cumplido voluntariamente lo dispuesto.

Juntas, estas dos unidades conforman el ciclo completo del proceso de ejecución dineraria en el derecho procesal argentino.

%===============================================================================
% CAPÍTULO 1: EL JUICIO EJECUTIVO (UNIDAD 22)
%===============================================================================

\chapter{El Juicio Ejecutivo}

\section{Ubicación en el Código Procesal}

El juicio ejecutivo se encuentra regulado en el Código Procesal Civil y Comercial de la Nación (en adelante CPCCN) a partir del artículo 520 y ss. Se lo localiza dentro de lo que el código denomina \textbf{proceso de ejecución}, junto con el tema de la ejecución de sentencias regulada en capítulos anteriores.

Sin embargo, es fundamental no confundir el juicio ejecutivo con el proceso de ejecución de sentencias, ya que ambos responden a una lógica completamente distinta. Mientras que la ejecución de sentencia tiene como antecedente un conocimiento pleno y exhaustivo del conflicto, el juicio ejecutivo opera sobre la base de un conocimiento sumario o acotado.

\section{Concepto y Naturaleza Jurídica}

\subsection{El proceso de conocimiento sumario}

El juicio ejecutivo constituye un \textbf{proceso de conocimiento sumario}, lo que quiere decir que el conocimiento del juez en relación al conflicto va a estar necesariamente acotado, disminuido o, como se ha señalado, seccionado.

El conocimiento del juez no va a extenderse más allá que el título cartular que se lleva para ejecutar. Esto lo diferencia de otros procesos sumarios que existen en las legislaciones provinciales, donde existe un conocimiento pleno de la materia litigiosa.

\subsection{La opción por la rapidez}

La razón de ser del juicio ejecutivo es la \textbf{rapidez}. El legislador ha buscado deliberadamente sacrificar la exhaustividad en el conocimiento en aras de la celeridad. El objetivo es que el acreedor que cuenta con ciertas prestaciones, ante el incumplimiento del deudor, tenga a su disposición una herramienta ágil para recobrar el crédito o las sumas de dinero adeudadas, sin necesidad de someterse a un juicio de conocimiento pleno.

En el juicio ejecutivo no existe:
\begin{itemize}
    \item Traslado de demanda con la amplitud del juicio ordinario.
    \item Excepciones con la extensión admitida en el procedimiento ordinario.
    \item Apertura de prueba sin restricciones.
    \item Alegatos de cierre.
    \item La regla de apelabilidad irrestricta de todas las decisiones.
\end{itemize}

Al contrario, el legislador ha establecido un proceso restringido, al cual sólo se puede acceder ante la presencia de un \textbf{título ejecutivo}.

\section{Diferencia con la Ejecución de Sentencias}

Es imprescindible establecer las diferencias fundamentales entre el juicio ejecutivo y la ejecución de sentencias:

\begin{table}[H]
\centering
\begin{tabularx}{\textwidth}{Xp{5cm}}
\toprule
\textbf{Ejecución de Sentencias} & \textbf{Juicio Ejecutivo} \\
\midrule
Requiere un juicio ordinario previo con conocimiento pleno & Sólo requiere un título ejecutivo \\
\midrule
La sentencia es el antecedente del trámite ejecutorio & El título cartular es el fundamento directo \\
\midrule
El deudor pudo discutir exhaustivamente la causa & La causa de la obligación no puede discutirse en el proceso ejecutivo \\
\midrule
La cosa juzgada es material para todas las cuestiones & La cosa juzgada es formal respecto de las cuestiones causales \\
\midrule
El embargo resultante es `ejecutorio' (de una sentencia firme) & El embargo resultante es `ejecutivo' (del título cartular) \\
\bottomrule
\end{tabularx}
\end{table}

\subsection{El desdoblamiento del conocimiento del conflicto}

El juicio ejecutivo realiza un \textbf{desdoblamiento} del conocimiento del conflicto: opta por un mecanismo procesal ágil mediante el cual, con una cantidad reducida de trámites y a partir de la existencia de un título que traiga aparejada fuerza ejecutiva, el acreedor puede hacerse más rápidamente del crédito.

Esta opción por la rapidez implica la restricción de las defensas que podrá oponer el deudor. Las excepciones sólo podrán referirse a \textbf{cuestiones que emerjan del título cartular} que se está ejecutando. Para resguardar adecuadamente el derecho de defensa del deudor, el código le concede la posibilidad de discutir todas las demás cuestiones en el juicio ordinario posterior, en el cual la discusión ya no estará acotada.

\subsection{La cosa juzgada mixta en el juicio ejecutivo}

La cosa juzgada que emana del juicio ejecutivo presenta una naturaleza mixta:

\begin{itemize}
    \item \textbf{Cosa juzgada formal} respecto de las cuestiones causales que no pudieron ser debatidas en el ejecutivo. Estas podrán ser discutidas en el ordinario posterior.
    
    \item \textbf{Cosa juzgada material} respecto de aquellas cuestiones que sí pudieron ser debatidas en el marco del juicio ejecutivo. Por ejemplo, la excepción de pago es una excepción admisible en el ejecutivo: si se ventila y se resuelve, esa cuestión no podrá volver a plantearse en el ordinario posterior.
\end{itemize}

\section{Requisitos para Acceder a la Vía Ejecutiva}

Para poder acceder a la vía ejecutiva se requieren \textbf{dos condiciones esenciales}:

\begin{important}
1. Contar con un \textbf{título que traiga aparejada la ejecución} (título ejecutivo con fuerza ejecutiva).

2. Que exista una \textbf{obligación exigible de dar cantidades líquidas de dinero} o fácilmente liquidables.
\end{important}

\subsection{La exigibilidad de la obligación}

La obligación debe tener un \textbf{plazo vencido}, es decir, que sea \textbf{exigible}, que el acreedor pueda pedir su cobro en el momento de iniciar la ejecución.

Debe tratarse de una \textbf{obligación de dar sumas de dinero}. El juicio ejecutivo, al menos en la forma actualmente regulada, queda limitado a este tipo de obligaciones, excluyendo las obligaciones de hacer o de no hacer.

\subsection{La liquidez o fácil liquidabilidad}

\textbf{Líquida} significa que la cantidad de dinero adeudada se encuentra expresada en números específicos.

\textbf{Fácilmente liquidable} significa que si existen, por ejemplo, cánones mensuales (como en el caso de alquileres), la suma líquida puede obtenerse mediante una liquidación sencilla en el marco del proceso ejecutivo mismo.

\subsection{El caso de la moneda extranjera}

Una situación común se presenta en hipotecas pactadas en moneda extranjera. En Argentina, a lo largo del tiempo, los dólares han recibido distintos tratamientos jurídicos: en algunos períodos se los consideraba como una obligación de dar cosas, en otros como equivalente a moneda corriente.

Actualmente, conforme jurisprudencia consolidada, se los trata como una \textbf{obligación de dar}. Por lo tanto, lo que corresponde hacer es convertir la moneda extranjera a su equivalente en moneda nacional, de manera que se obtenga esa suma líquida por la cual pueda ejecutarse conforme al artículo 520.

\section{Títulos Ejecutivos---Artículo 523 CPCCN}

El artículo 523 del CPCCN realiza una enumeración de los supuestos en los cuales se puede acceder a la vía ejecutiva. Es importante señalar que la vía ejecutiva es siempre \textbf{optativa} para el acreedor: este podría optar por un juicio ordinario directamente conforme al artículo 521.

\begin{formula}{Títulos con Fuerza Ejecutiva (art. 523 CPCCN)}

1°) El instrumento público presentado en forma.

2°) El instrumento privado suscripto por el obligado, reconocido judicialmente, o cuya firma estuviese certificada por escribano con intervención del obligado y registrada la certificación en el protocolo.

3°) La confesión de deuda líquida y exigible prestada ante el juez competente para conocer en la ejecución.

4°) La cuenta aprobada o reconocida como consecuencia del procedimiento establecido en el artículo 525.

5°) La letra de cambio, factura de crédito, cobranza bancaria de factura de crédito, vale o pagaré, el cheque y la constancia de saldo deudor en cuenta corriente bancaria, de acuerdo a las condiciones establecidas por las leyes en la materia.

6°) El crédito por alquileres o arrendamiento de muebles.

7°) Los demás títulos que tuvieran fuerza ejecutiva por ley y no estén sujetos a un procedimiento especial.
\end{formula}

\subsection{El inciso 7°: apertura a títulos especiales}

El inciso 7° abre la posibilidad de acceder a la vía ejecutiva a cualquier título que una ley especial le otorgue fuerza ejecutiva. Ejemplos de esto incluyen:

\begin{itemize}
    \item Deudas por matrículas en colegios profesionales (que cuentan con un certificado de deuda emitido por la entidad).
    \item Deudas fiscales con ciertos requisitos.
    \item Créditos por expensas en consorcios, que se conforman mediante certificado de deuda suscripto por el administrador del consorcio.
\end{itemize}

Estos títulos especiales tienen fuerza ejecutiva directamente de las leyes que los crean.

\section{Títulos Incompletos y Preparación de la Vía Ejecutiva}

Dentro de la enumeración de títulos ejecutivos se encuentran los denominados \textbf{títulos incompletos}: instrumentos que no tienen per se fuerza ejecutiva, sino que requieren algún trámite previo para lograr la preparación de la vía ejecutiva.

\subsection{Ejemplo: el contrato de locación}

Cuando el acreedor (locador) desea ejecutar un contrato de locación por el cobro de alquileres, no es posible determinar de inmediato cuál es la suma líquida a ejecutar. Ello se debe a que en el momento de la celebración del contrato no se sabía en qué mes el inquilino iba a dejar de pagar.

Por ello, existe un trámite previo: se cita al inquilino ante el tribunal para que acompañe el último recibo de pago del alquiler. A partir de ahí, todos los períodos posteriores sin recibo se consideran adeudados, y mediante una sencilla cuenta aritmética se obtiene la suma fácilmente liquidable.

Si el contrato de locación no tiene la firma del inquilino certificada por escribano, el inquilino deberá reconocer o desconocer la firma. La sanción por no concurrir es que se lo tendrá por reconocido. Si concurre y niega la firma, y luego se prueba mediante pericial caligráfica que le pertenece, le corresponde una multa del treinta por ciento más sobre la deuda reclamada.

\section{El Embargo Ejecutivo}

Cuando se tiene un título completo con fuerza ejecutiva, o un título incompleto que luego se preparó, el paso siguiente es directamente la intimación al deudor para que oponga excepciones, y la traba del embargo.

\begin{table}[H]
\centering
\begin{tabularx}{\textwidth}{>{\raggedright\arraybackslash}p{2.5cm}>{\raggedright\arraybackslash}p{3.5cm}>{\raggedright\arraybackslash}X}
\toprule
\textbf{Tipo de Embargo} & \textbf{Fundamento} & \textbf{Momento} \\
\midrule
Preventivo & Medida cautelar: requiere verosimilitud del derecho y peligro en la demora & Antes o durante el proceso, por solicitud de parte \\
\midrule
Ejecutivo & Fuerza ejecutiva del título cartular & Al inicio del juicio ejecutivo, sin necesidad de sentencia previa \\
\midrule
Ejecutorio & Sentencia firme que se ejecuta & En la etapa de ejecución de sentencia \\
\bottomrule
\end{tabularx}
\end{table}

El \textbf{embargo ejecutivo} es aquel que se traba a consecuencia de la fuerza ejecutiva de un título, que es lo que ocurre en el marco del juicio ejecutivo.

\section{La Intimación de Pago---Artículo 531 CPCCN}

La intimación de pago equivale en el juicio ejecutivo a lo que es el traslado de demanda en el juicio ordinario, pero se realiza mediante un mandamiento.

El mandamiento---que dice de \textbf{``intimación de pago y citación de remate''}---se libra al domicilio del deudor. Al deudor se le otorgan \textbf{cinco días hábiles judiciales} para interponer las excepciones que estime con derecho a poner.

El mandamiento puede ser acompañado del embargo domiciliario---en cuyo caso el oficial de justicia va al domicilio del deudor, lo intima de pago y embarga los muebles que sean embargables--- o ser simplemente un mandamiento de intimación de pago, dejando el embargo para otro carril: por ejemplo, si se trata de inmuebles, se libra un oficio directamente al Registro de la Propiedad Inmueble; si se habla de cuentas bancarias, a la entidad bancaria donde esté registrada esa cuenta.

\begin{important}
Conforme al artículo 531, son \textbf{irrenunciables} para el deudor:

1. La intimación de pago.
2. La citación para oponer excepciones.
3. La sentencia.
\end{important}

\section{Ampliación de la Ejecución---Artículos 540 y 541}

En el curso del juicio ejecutivo, especialmente en las ejecuciones de expensas, pueden surgir nuevos períodos de la misma obligación mientras tramita el juicio.

\begin{formula}{Ampliación de Ejecución}

\textbf{Artículo 540:} Ampliación por vencimientos anteriores a la sentencia. Se reitera la intimación de pago y los nuevos períodos pasan a integrar la suma líquida reclamada.

\textbf{Artículo 541:} Ampliación por vencimientos con posterioridad a la sentencia. Una vez firme la sentencia de trance y remate, sólo se permite que el deudor traiga documentación de pago de los períodos posteriores. No se abre nueva discusión.
\end{formula}

\section{Excepciones Admisibles---Artículo 544}

Una vez notificado el ejecutado, tiene cinco días para interponer las excepciones previstas y taxativamente enumeradas en el artículo 544 del CPCCN.

\begin{formula}{Excepciones Admisibles en el Juicio Ejecutivo (art. 544 CPCCN)}

1°) Incompetencia.

2°) Falta de personería en el ejecutante o en el ejecutado, por carecer de capacidad civil para estar en juicio.

3°) Litispendencia.

4°) Falsedad o inhabilidad del título con el que se pide la ejecución.

5°) Prescripción.

6°) Pago total o parcial documentado.

7°) Compensación de crédito líquido que resulte de documento que traiga aparejada la ejecución.

8°) Quita, espera, remisión, novación, transacción, conciliación o cosa juzgada.
\end{formula}

\begin{note}
La falta de legitimación ha sido omitida expresamente por el código, pero se ha incorporado jurisprudencial y doctrinariamente a través de la excepción de inhabilidad de título.
\end{note}

\section{Falsedad e Inhabilidad del Título (Art. 544, inc. 4°)}

\subsection{Falsedad material vs. falsedad ideológica}

Es fundamental establecer una distinción entre dos tipos de falsedad:

\begin{table}[H]
\centering
\begin{tabularx}{\textwidth}{Xp{5cm}}
\toprule
\textbf{Falsedad Material (Admisible)} & \textbf{Falsedad Ideológica (Inadmisible)} \\
\midrule
Adulteración física del documento & Cuestiona la causa de la obligación \\
\midrule
Ejemplo: tachaduras, sobreescrituras, fecha alterada, nombre agregado & Ejemplo: alegar que no se recibió el dinero del préstamo \\
\midrule
Se prueba con pericial caligráfica & Debe discutirse en juicio ordinario posterior \\
\midrule
Causa cosa juzgada material para ese punto & Produce sólo cosa juzgada formal en el ejecutivo \\
\bottomrule
\end{tabularx}
\end{table}

La \textbf{falsedad material} es la adulteración física del documento: tachaduras, sobreescrituras, alteración de fecha, agregación de nombres posteriores a la firma. Todos estos son supuestos de falsedad material que constituyen una excepción admisible en el juicio ejecutivo.

La \textbf{falsedad ideológica} es aquella que cuestiona la causa de la obligación. Por ejemplo, alegar que aunque se recibió un título representativo de un préstamo, en realidad no se recibió el dinero. Esta falsedad ideológica \textbf{no es admisible} en el marco del juicio ejecutivo: debe discutirse en el juicio ordinario posterior.

\subsection{Los casos extremos: ordinarización del proceso}

Existen casos extremos en los que se ha presentado una situación de coerción o amenaza que ha viciado el consentimiento del obligado. Por ejemplo, se obliga a una persona a firmar una hipoteca bajo amenazas, sin que reciba contraprestación alguna. Posteriormente, los acreedores (quien realizó la coerción) ejecutan esa hipoteca y subastan la casa.

En estos casos extremos, la jurisprudencia ha admitido la \textbf{ordinarización del proceso ejecutivo}, permitiendo que el juicio ordinario corra de manera paralela al ejecutivo. En el ordinario, el deudor podría pedir como medida cautelar la suspensión del juicio ejecutivo, o al menos que ambas sentencias se dicten de modo conjunto.

\subsection{La inhabilidad del título}

La excepción de \textbf{inhabilidad de título} permite cuestionar que el instrumento presentado no trae aparejada la ejecución, lo cual obligará al juez a examinar nuevamente el título en la sentencia.

También se utiliza para plantear la \textbf{falta de legitimación}: que el acreedor que promueve la ejecución no es la persona que surge del título, o que el ejecutado no es quien debe pagar.

\begin{important}
Conforme al artículo 544, inciso 4°: tanto la excepción de falsedad como la de inhabilidad del título sólo pueden interponerse siempre que el ejecutado \textbf{niegue la existencia de la deuda}.
\end{important}

\section{Apertura a Prueba y Sentencia de Trance y Remate}

Una vez presentadas las excepciones, el juez analiza si cumplen los requisitos legales. Las excepciones que no fueron autorizadas por ley son rechazadas conforme al artículo 547.

Si las excepciones son de puro derecho, el juez las resuelve sin abrir prueba. Si corresponde, puede abrirse prueba por un plazo breve determinado.

\textbf{La sentencia del juicio ejecutivo} se denomina \textbf{sentencia de trance y remate}, y su plazo de dictado es de \textbf{diez días}, ya que se la equipara a una resolución interlocutoria.

\begin{note}
El artículo 547 autoriza al juez a rechazar in limine (de plano) todas las excepciones que no se encuentren previstas legalmente.
\end{note}

\section{Apelación e Inapelabilidad en el Juicio Ejecutivo}

En el juicio ejecutivo, la regla es \textbf{contraria} a la del juicio ordinario: la \textbf{inapelabilidad es la norma}, y sólo son apelables aquellas decisiones que el código autoriza expresamente. Cualquier apelación que se deduzca fuera de los casos permitidos será denegada.

\begin{formula}{Casos de Apelación de la Sentencia de Trance y Remate (art. 554 CPCCN)}

Son apelables cuando:

a) Se desestiman sin más las excepciones por no cumplirse los requisitos del artículo 544.

b) Las excepciones hubiesen tramitado como de puro derecho.

c) Se hubiese producido prueba respecto de las opuestas.

d) La sentencia versare sobre puntos ajenos al ámbito natural del proceso o que cause gravamen irreparable, casos en los que se deja expedito el juicio ordinario posterior.

\textbf{Recurso aplicable:} Apelación en relación (más acotada que la apelación libre del juicio ordinario).
\end{formula}

\section{Costas y Juicio Ordinario Posterior}

Las costas en el juicio ejecutivo siguen la regla general: las soporta la parte vencida. Existe una salvedad cuando el ejecutado logra que se rechace \textbf{parcialmente} la ejecución---por ejemplo, acreditando el pago parcial documentado---. En esos casos no se aplica una exención proporcional de costas: la regla es que la ejecución prospera integralmente respecto de costas.

\begin{important}
Conforme al artículo 553: hasta que no esté concluido el juicio ejecutivo no se podría promover el juicio ordinario posterior, salvo en los casos de ordinarización paralela ya vistos.
\end{important}

\section{Ejecuciones Especiales---Artículo 595 y ss.}

A partir del artículo 595 el código regula las \textbf{ejecuciones especiales}: distintos supuestos de subastas de bienes muebles, inmuebles y otras cuestiones especializadas. Se prevén reglas específicas para algunos tipos de ejecución que se basan en el trámite del juicio ejecutivo, pero que presentan condiciones particulares.

\subsection{La ejecución hipotecaria}

Con la Ley 24.441 se creó lo que se denomina la \textbf{hipoteca privada}: un mecanismo de ejecución que se apoya parcialmente en la jurisdicción para ciertos trámites, pero que fundamentalmente se lleva adelante de modo privado.

Esta regulación ha sido objeto de críticas doctrinarias porque, entre otras cosas, omite el dictado de una sentencia cuando el deudor no opone excepciones. Esto presenta una objeción constitucional importante: la Constitución Nacional establece que nadie puede ser privado de su propiedad sino en virtud de una sentencia previa. Si el deudor no interpone excepciones, el juez no tiene obligación de dictar sentencia, violando así este principio.

%===============================================================================
% CAPÍTULO 2: CUMPLIMIENTO DE LA SENTENCIA DE REMATE (UNIDAD 23)
%===============================================================================

\chapter{Cumplimiento de la Sentencia de Remate}

\section{La Etapa Coercitiva del Proceso}

Hasta la sentencia definitiva del juicio ejecutivo, toda la discusión---la demanda, la contestación de excepciones, la posible apertura a prueba, la parte recursiva---tuvo como objetivo obtener un reconocimiento de derechos en sede jurisdiccional.

Ahora corresponde llevar a la \textbf{realidad práctica} lo que se dispuso en el proceso judicial. Los derechos reconocidos judicialmente deben tener su correlato en prestaciones concretas, conductas de las personas o sujetos en la vida de relación. El deudor debe pagar la deuda que fue reconocida judicialmente; debe hacer o dejar de hacer las actividades que la sentencia ordena.

Los elementos de la \textbf{jurisdicción} pueden enunciarse en tres:

\begin{enumerate}
    \item \textbf{Notio} (conocimiento del conflicto).
    \item \textbf{Documentio} (regulación del procedimiento).
    \item \textbf{Coercio} (coerción estatal: la imposición de la decisión por la fuerza cuando las personas no han cumplido voluntariamente).
\end{enumerate}

\section{Títulos Asimilables a la Sentencia}

Tanto el cumplimiento de la sentencia de trance y remate como el procedimiento de ejecución de sentencia---que se utiliza cuando el juicio ordinario está concluido---sirven también para aquellos títulos \textbf{asimilables a la sentencia}:

\begin{itemize}
    \item Acuerdos de mediación.
    \item Honorarios regulados judicialmente.
    \item Sentencias extranjeras reconocidas por el mecanismo del exequátur.
    \item Otros que las leyes especiales hayan asimilado.
\end{itemize}

\section{Inapelabilidad del Ejecutado---Artículo 560}

El ejecutado que llegó a esta etapa de cumplimiento de sentencia ya tuvo \textbf{amplia posibilidad de defensa} durante todo el proceso previo. Por ello, en el cumplimiento de la sentencia de remate, la regla es la \textbf{inapelabilidad para el ejecutado}.

\begin{formula}{Inapelabilidad en Cumplimiento de Sentencia (art. 560 CPCCN)}

Son inapelables para el ejecutado todas las resoluciones que se dicten durante el trámite de cumplimiento de la sentencia de remate.

Excepciones:

a) Resoluciones que no puedan constituir objeto del juicio ordinario posterior.

b) Aquellas que debiendo ser objeto del juicio ordinario posterior ya fueron debatidas.

c) Resoluciones que se relacionen al carácter de parte del ejecutado.

d) Los casos del artículo 550 (puntos ajenos al ámbito natural del proceso que causen gravamen irreparable).
\end{formula}

\section{El Embargo Ejecutorio: Presupuesto del Remate}

\begin{important}
No hay posibilidad de que el proceso avance hacia el remate de un bien si no lo ha embargado previamente.
\end{important}

Se está refiriendo al \textbf{embargo ejecutorio}, que es el resultado de la traba que se realiza una vez que hay una condena firme en sede jurisdiccional.

Si lo que se va a percibir son sumas de dinero ya embargadas en cuentas bancarias, no hay necesidad de realizar un remate: directamente con la firmeza de la sentencia se practica una liquidación y se ordena el giro de esos fondos al acreedor.

Lo relevante del remate aparece cuando hay un \textbf{bien específico}---mueble o inmueble---que necesita ser subastado para transformarse en dinero, que luego pueda aplicarse al pago del crédito.

\section{Bienes Subastables}

Se puede subastar \textbf{todo aquello que esté en el comercio}, es decir, todo aquello que pueda ser vendido conforme a la ley.

También es posible subastar \textbf{derechos hereditarios}. Cuanto mayor sea la determinación del acervo hereditario, con mayor precisión se va a poder establecer el valor de ese derecho y, de tal manera, fijar la base de la venta.

\section{El Martillero: Designación y Funciones}

El encargado de llevar adelante la subasta a las órdenes del juez es el \textbf{martillero o corredor inmobiliario}, que se anota en las listas de las cámaras profesionales para ser designado. En el fuero civil de la Nación, la designación actualmente se realiza directamente por sorteo por computadora.

\subsection{Obligaciones del martillero}

\begin{enumerate}
    \item \textbf{Cumplir las funciones personalmente:} no puede delegarlas a terceros.
    \item \textbf{Responder sólo a las órdenes del juez:} no de las partes interesadas.
    \item \textbf{Aceptar el cargo obligatoriamente} una vez sorteado.
    \item \textbf{Informar al juez} el valor fiscal y de mercado del bien a subastar.
    \item \textbf{Exhibir el inmueble} a los eventuales compradores en los días y horas fijados.
    \item \textbf{Publicar los edictos} anunciando la subasta.
    \item \textbf{Depositar en el proceso} las sumas percibidas en el acto de subasta.
    \item \textbf{Solicitar un anticipo de fondos} al ejecutante para cubrir los gastos del trámite (publicación, exhibición, etc.).
\end{enumerate}

\section{La Liga: Distorsión del Sistema de Subastas}

\begin{example}{Fenómeno de la ''Liga'' en Subastas Judiciales}

Una serie de personas se dedica a distorsionar lo que se obtiene en las subastas judiciales mediante un mecanismo conocido como la ``liga''.

\textbf{Mecanismo:} Impiden que vayan eventuales compradores interesados en un bien particular, de modo de intentar adquirirlo por el menor valor posible---si es posible, por la base de la subasta---. Luego hacen la ``verdadera subasta'' fuera del tribunal: el verdadero comprador le paga a la liga la diferencia entre lo que pagaron en sede judicial y el valor real del bien.

\textbf{Impacto económico:} La diferencia---que puede oscilar entre el 20 y el 30\% del valor de un inmueble---es dinero que el deudor no va a obtener si hay un remanente. El acreedor cobra menos que lo que podría haber cobrado, y el deudor pierde también.

\textbf{Ejemplo numérico:} Si un inmueble vale \$100.000 y la liga logra que se subaste a \$60.000 cuando podría haberse vendido a \$90.000, esos \$30.000 van al bolsillo de la liga.
\end{example}

\section{Subastas Electrónicas: Una Necesidad Impostergable}

Resulta anacrónico que en la actualidad la práctica judicial siga haciendo que se saque a subasta un inmueble en una oficina del tribunal, donde todos deben estar presentes durante aproximadamente dos horas para hacer la puja en forma presencial.

La \textbf{Provincia de Buenos Aires} tiene una ley de subasta electrónica que se está implementando gradualmente. En este sistema, la subasta no dura dos horas sino \textbf{quince días}: distintos interesados toman conocimiento por internet del inmueble, se registran en la plataforma y hacen ofertas que se actualizan a medida que aparece un postor con mejor oferta.

\textbf{Ventajas múltiples:}

\begin{itemize}
    \item Elimina a la liga y el fraude sistemático.
    \item Obtiene valores mucho más altos y reales.
    \item Beneficia tanto al acreedor---que tiene más dinero---como al deudor---que puede obtener un remanente luego de pagada la deuda.
    \item Amplía significativamente el universo de potenciales compradores.
\end{itemize}

\section{Trámites Previos a la Subasta de Inmuebles}

\subsection{Informes registrales y de deudas}

Cuando el juez ordena la subasta y designa al martillero, ordena también que las partes se ocupen de obtener informes de todos los entes sobre las deudas que tienen los inmuebles.

\begin{important}
Es obligatorio obtener informes de:

\begin{itemize}
    \item Municipalidades y organismos de recaudación provincial (como ARBA en Buenos Aires) sobre deudas impositivas.
    \item Organismos de servicios: agua (AySA), gas, luz, etc.
    \item Registro de la Propiedad Inmueble: sobre embargos, inhibiciones y gravámenes (hipotecas) que pesen sobre el inmueble y sobre el ejecutado.
\end{itemize}

\textbf{Validez de los informes:} 60 días. No es necesario actualizarlos antes de la subasta, salvo que haya transcurrido mucho tiempo por alguna suspensión del proceso.

Para inmuebles en propiedad horizontal: debe requerirse también la deuda por expensas.
\end{important}

\subsection{Citación de los acreedores con gravámenes}

No es lo mismo subastar un inmueble sin gravamen que uno con hipotecas y embargos previos. Cuando el inmueble tiene hipotecas, embargos u otros gravámenes, el acreedor ejecutante debe respetar las preferencias y privilegios de esos embargos en la distribución del producido.

Una vez que el inmueble se subasta y el embargo pasa al producido de la subasta por \textbf{subrogación real}, el juez tiene que repartir ese dinero de acuerdo a los privilegios y preferencias de los gravámenes y embargos.

Por eso se notifica a todos los acreedores con gravámenes \textbf{antes de la subasta}, para que puedan presentarse, defender el precio de venta y controlar la legalidad de todo el proceso.

\section{Los Edictos y la Publicidad---Artículo 566}

\begin{formula}{Publicidad de la Subasta (art. 566 CPCCN)}

El remate se anunciará por dos días en el Boletín Oficial y en otro diario de amplia circulación del lugar donde se ubica el inmueble.

\textbf{Contenido obligatorio del edicto:}

\begin{itemize}
    \item Expediente, carátula y juzgado donde tramita el proceso.
    \item Fecha, hora y lugar de la subasta.
    \item Días y horas en que el martillero exhibirá el inmueble a los interesados.
    \item Deudas del inmueble (servicios, tasas, contribuciones municipales, expensas).
    \item Gravámenes que pesan sobre el bien.
    \item Obligación del adquirente de pagar en el acto la seña y la comisión del martillero.
    \item Modalidades especiales fijadas por el juez para la venta.
\end{itemize}
\end{formula}

\begin{note}
Existe una crítica contemporánea sobre esta regulación: en la actualidad, el diario impreso como elemento de comunicación es anacrónico. Muy poca gente sigue comprando diarios. Sin embargo, a la fecha, la norma sigue vigente y debe cumplirse. El martillero puede solicitar al tribunal autorización para realizar publicidad adicional---por ejemplo, por internet---para atraer más compradores.
\end{note}

\section{La Base de la Subasta---Artículo 578}

El martillero tiene la obligación de informar al juez el \textbf{valor fiscal} de ese inmueble y el \textbf{valor de mercado}. Con frecuencia, los valores fiscales se encuentran muy desconectados del verdadero valor de mercado de una propiedad.

\begin{formula}{Fijación de la Base de Venta (art. 578 CPCCN)}

Si no hay acuerdo de partes, la base de venta se fijará en \textbf{dos tercios de la valuación fiscal actualizada}.

\textbf{Práctica judicial:} En general el juez se aparta de ese valor fiscal para impedir que los bienes sean malvendidos. El juez fija un valor menor al de mercado puro con el objetivo de atraer distintos compradores y lograr una mejor venta.

Las partes---y los acreedores con gravámenes---pueden impugnar la tasación informada por el martillero.
\end{formula}

\section{El Acto de Remate}

Hay distintos resultados posibles en el acto de remate: que haya un comprador interesado, que el remate fracase por falta de postores, o que el propio acreedor ejecutante quiera adquirir el bien.

\subsection{Cuando hay comprador}

Una vez que se realiza el acto de remate y hay un comprador interesado, el martillero le hace firmar al comprador el \textbf{boleto de venta judicial}, donde consta:

\begin{itemize}
    \item El precio de compra acordado.
    \item Lo que paga de seña (generalmente el 30\% del precio).
    \item Lo que paga de comisión del martillero.
    \item Otros impuestos que correspondan a esa venta.
\end{itemize}

El martillero retiene su comisión y deposita el dinero obtenido en el proceso, informando el nombre e identidad del comprador.

El comprador debe constituir domicilio electrónico conforme al artículo 41 CPCCN para ser notificado de las decisiones judiciales posteriores.

\begin{formula}{Secuencia Tras el Acto de Remate}

1) El martillero informa el resultado de la subasta al tribunal.

2) Se da conocimiento del resultado al acreedor y al deudor: tienen \textbf{5 días} para formular objeciones a la venta.

3) Si no hay objeciones o son desestimadas: el juez aprueba formalmente el remate.

4) El comprador tiene \textbf{5 días} para depositar el saldo de precio (el 70\% restante, dado que el 30\% se paga como seña en el acto).
\end{formula}

\subsection{Crítica a la exigencia del dinero en efectivo}

La exigencia de que el comprador presente el 30\% en efectivo el día de la subasta presenta objeciones prácticas considerables en la actualidad. Requiere que las personas concurran con dinero que quizá no utilicen regularmente, con todos los riesgos que eso trae en términos de seguridad personal. Para los fines de la efectividad de la subasta es un requisito bastante complejo.

\section{Aprobación y Nulidades}

Los plazos para plantear nulidades son de \textbf{cinco días} contados desde que se conoce el resultado de la subasta. Para quien estuvo presente en el acto, el cómputo se inicia desde el propio acto, independientemente de cuándo el martillero informe al tribunal.

\begin{important}
Las nulidades que se pueden plantear en esta etapa son las nulidades del acto mismo: aquellas que ocurrieron el día de la subasta. No se pueden replantear las nulidades procesales anteriores (la fijación de la base, la publicación de edictos, etc.), porque cada una de ellas tenía su propio plazo de cinco días desde el momento en que se conocían en el expediente para ser cuestionadas.

En general, la única nulidad que se va a poder plantear es algún incumplimiento grave en el acto propio de la subasta.
\end{important}

\subsection{El principio de trascendencia}

Quien articula la nulidad tiene la obligación de \textbf{acreditar el perjuicio}. Sin un perjuicio concreto demostrado, no procede la nulidad. La falta menor no hace caer el acto.

La trascendencia de la nulidad tiene que tener una entidad suficiente para hacer caer la venta como acto válido.

\section{El Sobreseimiento del Juicio---Artículo 583}

Existe aún después de realizada la subasta una última chance de que el deudor recupere el bien. El ejecutado podría, \textbf{antes de que el comprador en subasta deposite el saldo de precio}, pedir el sobreseimiento del juicio.

Una vez que el comprador ha depositado el saldo de precio, se considera que se \textbf{perfeccionó la venta} y ya no es posible el sobreseimiento.

\begin{formula}{Requisitos para el Sobreseimiento (art. 583 CPCCN)}

Para que proceda el sobreseimiento, el deudor debe depositar:

1) El capital adeudado.

2) Los intereses que correspondan según lo presupuestado.

3) Las costas del proceso.

4) Todos los gastos del procedimiento de subasta.

5) Un importe a favor del comprador integrado por:
   \begin{itemize}
       \item La comisión que le pagó al martillero.
       \item Los sellados del boleto de venta.
       \item El equivalente a la seña depositada.
   \end{itemize}
\end{formula}

\section{Título del Adquirente e Inscripción Registral}

Así como en una compraventa de un inmueble fuera de los tribunales hay una escritura pública ante escribano, en el marco de un proceso judicial lo que se otorga es un \textbf{testimonio}. Este testimonio contiene toda la referencia y los antecedentes de esa subasta judicial, y las distintas resoluciones del juez que permitieron que el adquirente pueda inscribir ese derecho.

El testimonio se extiende a través de un trámite más del proceso, que lo realiza el secretario del tribunal, sin necesidad de intervención de escribano público.

Este testimonio constituye el \textbf{nuevo título del bien} para el adquirente en subasta.

\subsection{Levantamiento de gravámenes}

El adquirente va a pedir al juez el levantamiento de todas las medidas dispuestas contra el inmueble---hipotecas, embargos, inhibiciones---al solo efecto de que se anote el nuevo título en el Registro de la Propiedad Inmueble.

Los acreedores que tenían esos gravámenes sobre el bien deben presentarse en el proceso a hacer valer sus preferencias y privilegios en la distribución del producido obtenido.

\section{Distribución del Producido---Liquidación}

Una vez depositado el saldo de precio por el comprador, se practica la \textbf{liquidación}: se computa el capital, los intereses, las costas y los gastos del proceso.

Se pone la liquidación en conocimiento del deudor ---quien tiene un plazo para impugnarla---y una vez aprobada se entrega el dinero al ejecutante, cancelando así el crédito.

\subsection{El remanente}

El deudor puede tener derecho a un \textbf{remanente} de la subasta: si el producido obtenido supera el total de la deuda (capital + intereses + costas + gastos), el excedente le corresponde al deudor.

\section{Subasta Fracasada y Acreedor Comprador}

\subsection{Fracaso por falta de postores}

Si en el acto de remate no hay postores interesados en comprar el inmueble, se realizará una \textbf{nueva subasta} conforme a los artículos 592 y 593 del CPCCN. En general, se intenta bajar la base de la nueva subasta para atraer nuevos interesados.

\subsection{El acreedor ejecutante como comprador}

El \textbf{acreedor ejecutante puede ir a comprar él mismo} el inmueble que se subasta, y \textbf{compensar su crédito} con el precio de compra.

En general, los acreedores ejecutantes que resultan adquirentes le piden al juez que los \textima de depositar la seña. Este constituye un gran beneficio para ese comprador: no tiene que concurrir con dinero en efectivo expuesto a riesgos de seguridad.

Esta misma posibilidad podría ser solicitada por algún acreedor embargante en otro juicio que tiene que venir a la subasta del inmueble.

\section{Los Plenarios: Un Aspecto Crítico para la Práctica}

\begin{important}
Diferencia fundamental según el fuero:

\textbf{Plenario Cámara Comercial:} El adquirente en subasta debe pagar todas las deudas que están expresadas en el edicto.

\textbf{Plenario Cámara Civil:} Todas las deudas del inmueble no las tiene que afrontar el adquirente en tanto alcance el dinero que se obtiene en la subasta. Es decir, si el producto de la subasta es insuficiente, el adquirente no responde personalmente por esas deudas.

\textbf{Consecuencia práctica:} Según si se subasta un inmueble en juzgado civil o comercial de la Nación, la ecuación económica para el potencial comprador cambia radicalmente.

\textbf{Recomendación:} Un abogado asesor debe verificar los plenarios vigentes en la jurisdicción antes de aconsejar a su cliente como potencial comprador de un inmueble en subasta.
\end{important}

\section{Desocupación del Inmueble---Artículo 589}

\begin{formula}{Momento de la Desocupación (art. 589 CPCCN)}

Como regla general: no hay desalojo o desocupación de los ocupantes del inmueble hasta que no se hubiese pagado el saldo de precio y realizada la tradición de la cosa al adquirente.

Excepciones: salvo que se haya pactado otra cosa específicamente. Por ejemplo, en los casos de hipotecas donde el hipotecado generalmente se obliga contractualmente a desocupar el inmueble dos días después del auto de subasta.
\end{formula}

%===============================================================================
% MÉTODO CORNELL
%===============================================================================

\chapter*{Método Cornell: Síntesis y Recuperación Activa}

\begin{table}[H]
\centering
\begin{tabularx}{\textwidth}{
    >{\raggedright\arraybackslash}p{0.25\textwidth}
    >{\raggedright\arraybackslash}X
}
\toprule
\textbf{Preguntas / Conceptos Clave} & \textbf{Notas / Respuestas / Explicación} \\
\midrule

\textbf{¿Qué es el juicio ejecutivo y cómo se ubica en el código?} & Es un proceso de conocimiento sumario (acotado) regulado a partir del art. 520 CPCCN. Prioriza la rapidez sobre la exhaustividad: el acreedor accede directamente al embargo y cobro sin un juicio ordinario previo. \\

\textbf{¿Por qué se lo denomina ``conocimiento sumario''?} & Porque el conocimiento del juez queda limitado al título cartular presentado, no se extiende a cuestiones causales que pueden debatirse luego en el ordinario posterior. \\

\textbf{¿Cuáles son los requisitos para acceder a la vía ejecutiva?} & 1) Tener un título ejecutivo (completo o preparado). 2) Obligación de dar sumas de dinero. 3) Que sea líquida o fácilmente liquidable. 4) Que sea exigible (plazo vencido). \\

\textbf{¿Qué diferencia hay entre títulos completos e incompletos?} & Completos traen aparejada la ejecución por sí solos (pagaré, cheque, escritura pública). Incompletos requieren un trámite previo de preparación (ej.: contrato de locación). \\

\textbf{¿Cuáles son los tipos de embargo?} & Preventivo: medida cautelar. Ejecutivo: en el juicio ejecutivo, por fuerza del título, sin sentencia previa. Ejecutorio: resultado de sentencia firme en ejecución. \\

\textbf{¿Qué excepciones admite el juicio ejecutivo?} & Taxativamente enumeradas en art. 544: incompetencia, falta de personería, litispendencia, falsedad o inhabilidad del título, prescripción, pago documentado, compensación, quita/espera/remisión/novación/transacción/conciliación/cosa juzgada. \\

\textbf{Falsedad material vs. falsedad ideológica: ¿cuál es admisible?} & Material (sí): adulteración física (tachaduras, sobreescrituras, fecha alterada). Ideológica (no): cuestiona la causa (``no recibí el dinero''). Esta última va al ordinario posterior. \\

\textbf{¿Qué es la cosa juzgada formal vs. material en el ejecutivo?} & Formal: cuestiones causales no debatibles en el ejecutivo pueden discutirse en el ordinario. Material: cuestiones debatidas en el ejecutivo (ej.: pago) generan cosa juzgada material. \\

\textbf{¿Cómo funciona la apelación en el juicio ejecutivo?} & Regla: inapelabilidad (inversa al ordinario). Solo son apelables los casos que autoriza el art. 554. Recurso: apelación en relación (más acotada). \\

\textbf{¿Qué es el cumplimiento de la sentencia de remate?} & Es la etapa coercitiva: una vez firme la sentencia, el Estado impone coercitivamente su decisión ejecutando bienes del deudor. Presupuesto: embargo ejecutorio. \\

\textbf{¿Cuáles son las funciones principales del martillero?} & Evaluar el bien, informar valor fiscal y de mercado, exhibirlo a compradores, publicar edictos, conducir el acto de remate, suscribir boleto, retener comisión y depositar producido. Solo responde al juez. \\

\textbf{¿Qué es ``la liga'' y cómo opera?} & Grupo que impide concurrencia de compradores para adquirir a precio mínimo, luego revende privadamente la diferencia (20-30\%). Perjudica a acreedor y deudor. \\

\textbf{¿Qué es una subasta electrónica y qué ventajas tiene?} & Sistema por internet durante 15 días donde interesados se registran y hacen ofertas. Elimina la liga, obtiene valores mayores, beneficia a acreedor y deudor. \\

\textbf{¿Cuál es la base de la subasta según art. 578?} & Si no hay acuerdo: dos tercios de valuación fiscal actualizada. El juez suele apartarse para fijar valor menor al de mercado y atraer compradores. \\

\textbf{¿Qué pasos siguen tras el acto de remate?} & 1) Martillero informa. 2) 5 días para objeciones. 3) Aprobación. 4) 5 días para depositar saldo. 5) Sobreseimiento antes del depósito. 6) Liquidación. 7) Testimonio e inscripción. \\

\textbf{¿Qué es el sobreseimiento (art. 583)?} & El ejecutado recupera el bien antes de que comprador deposite saldo. Debe pagar: capital + intereses + costas + gastos + comisión martillero + sellados + seña. \\

\textbf{¿Por qué importan los plenarios sobre deudas en subastas?} & Cámara Comercial: adquirente paga deudas en edicto. Cámara Civil: deudas se abonan con producido (no responde adquirente si insuficiente). Cambia ecuación económica. \\

\midrule

\multicolumn{2}{p{\textwidth}}{

\textbf{SÍNTESIS FINAL:} 

El juicio ejecutivo y el cumplimiento de la sentencia de remate constituyen los dos actos sucesivos del proceso de ejecución dineraria en Argentina. El juicio ejecutivo (Unidad 22) sacrifica exhaustividad por celeridad: permite agredir patrimonio sin juicio ordinario previo mediante un título ejecutivo. Las excepciones son taxativas (art. 544) y excluyen discusión de la causa, que se reserva al ordinario posterior. La cosa juzgada es mixta: formal para cuestiones causales, material para lo debatido. El cumplimiento de la sentencia de remate (Unidad 23) despliega poder coercitivo pleno: embargo ejecutorio, designación de martillero, informes registrales, publicidad mediante edictos, notificación a acreedores concurrentes, acto de remate. El producido se distribuye según privilegios de gravámenes. El adquirente obtiene título mediante testimonio del secretario sin escribana. Deficiencias: persistencia del sistema presencial que alimenta la ``liga'', y divergencia entre plenarios civiles y comerciales sobre deudas que asume el adquirente---variable que el operador jurídico debe verificar para asesoramiento económicamente preciso.

} \\

\bottomrule
\end{tabularx}
\end{table}

%===============================================================================
% BIBLIOGRAFÍA
%===============================================================================

\chapter*{Bibliografía}

\begin{itemize}
    \item Código Procesal Civil y Comercial de la Nación Argentina (CPCCN).
    \item Ley 24.441: Sobre Hipotecas.
    \item Artículos 520 a 595 del CPCCN: Regulación del Proceso de Ejecución y Juicio Ejecutivo.
    \item Legislación Provincial (Provincia de Buenos Aires): Disposiciones sobre Subastas Electrónicas.
\end{itemize}

%===============================================================================
% FIN DEL DOCUMENTO
%===============================================================================

\end{document}
